\documentclass[10pt,a4paper]{amsart}
\usepackage[margin=19mm]{geometry}
\usepackage{amsmath,amssymb,mathtools}
\usepackage[scr=rsfs,cal=euler]{mathalfa}
\usepackage{xfrac}
\usepackage[unicode,hidelinks,bookmarks=false]{hyperref}
\newcommand{\ie}{i.e.\ }
\newcommand{\cf}{cf.\ }
\newcommand{\resp}{respectively\ }
\newcommand{\Subsection}{\subsection}
\providecommand{\nobreakdash}{\nobreak\hbox{-}}
\setlength{\emergencystretch}{2em}
\multlinegap0pt
\usepackage{bm}
\usepackage{bbm}
\usepackage{dsfont}
\usepackage{xspace}
\usepackage{cancel}
\usepackage{mathtools}
\usepackage{amsthm}
\makeatletter
\@ifundefined{lemma}{\newtheorem{lemma}{Lemma}}{}
\@ifundefined{theorem}{\newtheorem{theorem}[lemma]{Theorem}}{}
\@ifundefined{definition}{\newtheorem{definition}[lemma]{Definition}}{}
\makeatother
\title[Bilinearity Preserving Algebraic Flux Correction]{Bilinearity Preserving Algebraic Flux Correction}
\author{Sarthak Sourav Dash, Abhinav Jha}
\date{Source: arXiv:2609.16830v1 (CC BY 4.0)}
\begin{document}
\maketitle

\noindent\emph{Source and licence.} Bilinearity Preserving Algebraic Flux Correction. Version arXiv:2609.16830v1 (CC BY 4.0), distributed under the Creative Commons Attribution 4.0 International licence, as recorded in the arXiv licence metadata for this exact version and shown on the abstract page https://arxiv.org/abs/2609.16830v1. Full rights notice: licenses/arxiv-2609.16830-CC-BY-4.0.txt.\par\smallskip

\noindent\textit{Editorial scope.} Counted unit: the Theorem whose source label is \texttt{thm:BMPiffBC} in the article (source0.tex lines 1086--1092, source label \texttt{thm:BMPiffBC}), delivered together with the article's own complete original proof of it (source0.tex lines 1094--1145, one \texttt{proof} environment of 52 lines). No mathematics is written, repaired or re-derived: statement and proof are the article's own text line for line. Required context retained uncounted: eq:BMP (lines 963--969); Def:BilMax (lines 965--967); eq:BMP2 (lines 1011--1017); lem:BMP2 (lines 1013--1015). Every definition, display and label of those lines is retained unchanged. Omitted from this package and not redistributed: 1--962, 968--1010, 1016--1085, 1093--1093, 1146--2810. Nothing inside a retained excerpt is omitted or altered: every retained body is delivered exactly as the article's lines are written, so no omission receipt of any kind accompanies this package. Citations: the retained excerpts carry no citation command at all, so no bibliography entry of the article is transcribed and no reference is redistributed. Reference adaptation: none was needed and none was made; every reference inside the delivered unit resolves inside it, so no unresolved reference or citation is produced. Printed slips kept literal: no printed word, symbol, hypothesis or number of the counted statement or of its proof was altered. Layout and class: the article's own class is \texttt{article} with packages including a4, amsmath, amssymb, amsthm, mathtools, graphicx, color, float, xspace, siunitx, enumitem, booktabs, tikz, pgfplots; the wrapper is the lane's pinned \texttt{amsart} editorial wrapper at 10pt on A4, which restates the theorem-like environments the retained text needs and adds the article's own macro definitions that the retained text needs (none), with extra wrapper packages (bm, bbm, dsfont, xspace, cancel, mathtools). The delivered numbering is the unit's own. Assets: no retained span contains a figure environment, a graphics-inclusion command, a PDF-page inclusion, a table or a TikZ picture; no image file is copied into this package, the package declares no asset dependency and no image file is redistributed with it. Licence: version arXiv:2609.16830v1 of this article is distributed under the Creative Commons Attribution 4.0 International licence, as recorded in the arXiv licence metadata for this exact version and shown on the abstract page https://arxiv.org/abs/2609.16830v1. A full rights notice is delivered at licenses/arxiv-2609.16830-CC-BY-4.0.txt. Characters: every retained excerpt is pure ASCII, so the delivered engine, XeLaTeX with its default Latin Modern text font, reports no missing character.
\begin{definition}
We will say a mesh patch $\omega_i$ satisfies \textbf{bilinear maximum principle} if 
\begin{equation}
p_i^{\min} \leq p_i \leq p_i^{\max} \qquad \forall p \in \mathbb{Q}_1(\omega_i), \label{eq:BMP}
\end{equation}
where $p_i^{\min} = \min\limits_{j \in S_i} p_j$ and $p_i^{\max}=\max\limits_{j \in S_i} p_j$, and $p_i$ is the value of $p$ at the $i$th node. Moreover equality in either inequality holds if and only if $p$ is constant.  A mesh will be said to satisfy bilinear maximum principle if and only if each mesh patch $\omega_i$ satisfies Eq.~\eqref{eq:BMP} corresponding to $i=1,2,\cdots,m$.	\label{Def:BilMax}
 \end{definition}

\begin{lemma}
A mesh satisfies Bilinear Maximum Principle if and only if it satisfies
\begin{equation}
p_i^{\min} \leq 0 \leq p_i^{\max} \qquad \forall p \in \mathbb{Q}_1(\mathbb{R}^2),\; \text{and} \; p_i =0.   \label{eq:BMP2}
\end{equation}
with $p_i^{\min} = \min\limits_{j \in S_i} p_j$, $p_i^{\max}=\max\limits_{j \in S_i} p_j$.\label{lem:BMP2}
\end{lemma}

\begin{theorem}
A mesh patch $\omega_i$ satisfies
\begin{equation*}
	p_i^{\min} \le p_i \le p_i^{\max} \qquad \forall\, p \in \mathbb{Q}_1(\omega_i)
\end{equation*}
if and only if it is bilinearly convex, and satisfies the bilinear maximum principle of Definition~\ref{Def:BilMax}, with equality only for constant $p$, 	if and only if it is strictly bilinearly convex.\label{thm:BMPiffBC}
\end{theorem}

\begin{proof}
Without loss of generality, let the coordinate system be shifted such that the central node $i$ is located at the origin, therefore $(x_i, y_i) = (0,0)$. Under this translation, any general bilinear function $p \in \mathbb{Q}_1(\mathbb{R}^2)$ that vanishes at the internal node ($p_i = 0$) takes the form:
$$
p(x,y) = bx + cy + dxy
$$
Let $ \mathbf{P_j} = (x_j, y_j, x_j y_j)^\top \in \mathbb{R}^3$ denote the lifted coordinates of the neighboring nodes $j \in S_i$, and let $\mathcal{C} = \text{conv}(\{\mathbf{P_j}\}_{j \in S_i})$ be the convex hull of these points. Since $S_i$ is a finite index set of neighboring nodes, $\mathcal{C}$ is a compact (closed and bounded) convex set in $\mathbb{R}^3$.

\textbf{($\impliedby$):} 
Suppose the mesh is bilinearly convex. Then the origin $\mathbf{0} = (0,0,0)^\top \in \mathcal{C}$ by definition. Algebraically, this is equivalent to the existence of a set of convex coefficients $\{\alpha_j\}_{j \in S_i}$ such that $\alpha_j \geq 0$, $\sum_{j \in S_i} \alpha_j = 1$, and
$$
\sum_{j \in S_i} \alpha_j \mathbf{P_j} = \mathbf{0}.
$$
Taking the inner product with an arbitrary $\mathbf{v} = (b,c,d)^\top \in \mathbb{R}^3$ (representing coefficients), both sides of the equality gives
$$
\mathbf{v} \cdot \sum_{j \in S_i} \alpha_j \mathbf{P_j} = \sum_{j \in S_i} \alpha_j (\mathbf{v} \cdot \mathbf{P_j}) = \sum_{j \in S_i} \alpha_j p(x_j, y_j) = 0
$$
	Since $\alpha_j \geq 0$ and adds up to $1$, the function values at the nodes $p(x_j, y_j)$ cannot all be strictly positive or negative unless they are all zero. Consequently, there exists at least one $j_+ \in S_i$ with  $p_{j_+}\geq 0$ and at least one $j_- \in S_i$ with  $p_{j_-}\leq 0$. This implies $p_i^{\min} \leq 0 \leq p_i^{\max}$, ( i.e. Eq.~\eqref{eq:BMP2}) which in turn, by Lemma~\ref{lem:BMP2},  is equivalent to Eq.~\eqref{eq:BMP}.
	
\textbf{ ($\implies$):} 
We will prove the converse by contraposition. Suppose the mesh patch is \textit{not} bilinearly convex. This implies  $\mathbf{0} \notin \mathcal{C}$.
	
Since $\mathcal{C}$ is a closed convex set and disjoint from  $\{\mathbf{0}\}$ (a closed convex set) , by the Hyperplane Separation Theorem, there exists a hyperplane strictly separating them. Which implies there exists a non-zero vector $\mathbf{v} = (b, c, d)^\top \in \mathbb{R}^3$ and $\epsilon > 0$ such that:
$$
\mathbf{v} \cdot \mathbf{0} < \epsilon < \mathbf{v} \cdot \mathbf{p} \quad \forall \mathbf{p} \in \mathcal{C}.
$$

Evaluating this at the extreme points $\mathbf{P_j} \in \mathcal{C}$ for all $j \in S_i$ gives:
\begin{equation}
\mathbf{v} \cdot \mathbf{P_j} > 0 \quad \forall j \in S_i. \label{eq:separation}
\end{equation}
Now, we construct a bilinear function $p \in \mathbb{Q}_1(\mathbb{R}^2)$ corresponding to the vector $\mathbf{v}$ :
$$
p(x,y) = bx + cy + dxy.
$$
By construction, this function vanishes at the central node ($p_i = p(0,0) = 0$). Evaluating $p$ at any neighboring node $j \in S_i$ yields its value as a simple dot product:
$$
p(x_j, y_j) = bx_j + cy_j + dx_j y_j = \mathbf{v} \cdot \mathbf{P_j}.
$$
Applying the strict separation inequality Eq.~\eqref{eq:separation}, we find that:
$$
p(x_j, y_j) > 0 \quad \forall j \in S_i.
$$

Because the function value is strictly positive at every neighboring node in the patch, the minimum value over the neighborhood must also be strictly positive:
$$
p_i^{\min} = \min_{j \in S_i} p(x_j, y_j) > 0.
$$

This contradicts Eq.~\eqref{eq:BMP2} and consequently condition Eq.~\eqref{eq:BMP} by Lemma~\ref{lem:BMP2} . Thus, if a mesh satisfies condition Eq.~\eqref{eq:BMP2}, it is bilinearly convex.

It remains to identify the strict case. If the patch is strictly bilinearly convex, let $p \in \mathbb{Q}_1(\omega_i)$ be non-constant with $p_i = 0$. Then $p$ does not vanish at every node of $S_i$, and since $\sum_{j \in S_i} \alpha_j p_j = 0$ with all $\alpha_j > 0$, the values $p_j$ cannot all be of one sign, so $p_i^{\min} < 0 < p_i^{\max}$. Conversely, if the origin lies in $\mathcal{C}$ but not in its relative interior, there is a supporting hyperplane, that is, a non-zero $\mathbf{v} = (b,c,d)^\top$ with $\mathbf{v} \cdot \mathbf{P_j} \geq 0$ for every $j \in S_i$ and $\mathbf{v} \cdot \mathbf{P}_{j_0} > 0$ for at least one $j_0 \in S_i$. The associated $p(x,y) = bx + cy + dxy$ is non-constant, vanishes at the internal node, and satisfies $p_i^{\min} = 0 = p_i$, so the bilinear maximum principle fails in its strict form.
\end{proof}

\end{document}
